\documentclass[12pt,a4paper]{article}

\usepackage[utf8]{inputenc}
\usepackage[T1]{fontenc}
\usepackage[romanian]{babel}
\usepackage{geometry}
\usepackage{booktabs}
\usepackage{array}
\usepackage{longtable}
\usepackage{amsmath}
\usepackage{hyperref}
\usepackage{xcolor}
\usepackage{listings}

\geometry{margin=2.5cm}

\definecolor{codebackground}{rgb}{0.95,0.95,0.95}
\lstdefinestyle{codestyle}{
    backgroundcolor=\color{codebackground},
    basicstyle=\ttfamily\small,
    breaklines=true,
    frame=single,
    keepspaces=true,
    showstringspaces=false,
    tabsize=2
}
\lstset{style=codestyle}

\title{Evolutia si imbunatatirea proiectului de Speech Enhancement}
\author{}
\date{}

\begin{document}

\maketitle
\tableofcontents
\newpage

\section{Introducere}

Proiectul are ca obiectiv reducerea zgomotului din semnale de vorbire printr-o
combinatie de metode clasice de procesare digitala a semnalului si metode bazate
pe invatare profunda. In forma actuala, sistemul include algoritmi DSP de
referinta, un detector VAD, modele neuronale de tip MaskNet, mecanisme de
evaluare, preseturi experimentale, un demo offline si o interfata de demonstratie
in timp real.

Documentul prezinta, pas cu pas, imbunatatirile aduse proiectului fata de primele
iteratii. Comparatia este realizata pe baza structurii actuale a proiectului, a
istoricului Git si a componentelor pastrate pentru compatibilitate, precum
\texttt{SimpleMaskNet}. Primele iteratii pot fi caracterizate ca o varianta
functionala, dar sumara: accentul era pus pe ideea generala de denoising cu
MaskNet, metode clasice si VAD, fara un pipeline experimental complet standardizat.
Forma actuala este o platforma mai matura, modulara si reproductibila.

\section{Imagine de ansamblu asupra evolutiei}

In primele iteratii, proiectul era descris sintetic ca un sistem de imbunatatire a
semnalului audio, cu MaskNet, metode DSP, VAD, evaluare si demo. Ulterior,
proiectul a fost extins in mai multe directii:

\begin{enumerate}
    \item Organizarea codului a fost separata pe module clare: date, DSP, modele,
    antrenare, evaluare, utilitare si demo in timp real.
    \item Modelul neuronal a evoluat de la o retea convolutionala simpla la o
    arhitectura U-Net cu skip connections, BiLSTM, attention si blocuri reziduale.
    \item Pipeline-ul de date a fost extins cu perechi clean/noisy, STFT calculat
    on-the-fly, padding pentru secvente variabile, resampling la 16 kHz si suport
    pentru etichete VAD hard/soft.
    \item Antrenarea a devenit configurabila prin preseturi pentru model,
    optimizer, scheduler, loss si dataset.
    \item Evaluarea a trecut de la validari punctuale la comparatii intre metode,
    checkpoint-uri si seturi de etichete VAD.
    \item Proiectul a primit instrumente de livrare: CLI unificat, healthcheck,
    logging, inventariere de artefacte si clasificarea rezultatelor oficiale.
    \item Demo-ul a evoluat de la rulare offline la monitorizare in timp real,
    cu selectie de metoda, vizualizari, fisiere externe, scenarii si inregistrare.
    \item In cea mai recenta etapa a fost adaugata varianta Complex MaskNet, care
    estimeaza o masca complexa, nu doar o masca reala pe magnitudine.
\end{enumerate}

\section{Comparatie generala intre primele iteratii si forma actuala}

\begin{longtable}{p{0.23\textwidth}p{0.34\textwidth}p{0.34\textwidth}}
\toprule
\textbf{Aspect} & \textbf{Primele iteratii} & \textbf{Forma actuala} \\
\midrule
\endhead
Structura proiectului &
Structura mai redusa, centrata pe functionalitatea principala. &
Structura modulara in \texttt{src/data\_prep}, \texttt{src/dsp},
\texttt{src/models}, \texttt{src/train}, \texttt{src/eval}, \texttt{src/tools}
si \texttt{src/utils}. \\
\addlinespace
Model neuronal &
Model simplu de estimare a mastii, reprezentat prin \texttt{SimpleMaskNet}. &
MaskNet avansat cu U-Net, skip connections, BiLSTM, channel attention, variante
reziduale, variante rapide si suport pentru masca complexa. \\
\addlinespace
Date &
Utilizare de perechi audio clean/noisy intr-o forma de baza. &
Suport pentru VoiceBank-DEMAND si LibriSpeech+DEMAND, matching robust de fisiere,
resampling, crop, augmentare, preprocessing si VAD labels. \\
\addlinespace
Antrenare &
Configuratie mai putin standardizata si mai greu de reprodus. &
Preseturi oficiale, Adam/AdamW, OneCycleLR, mixed precision, checkpoint-uri,
resume, istoric de antrenare si sumar JSON. \\
\addlinespace
Evaluare &
Evaluare punctuala a unor metode. &
Evaluare comparativa intre noisy baseline, Spectral Subtraction, Wiener Filter si
MaskNet, plus comparare de checkpoint-uri, VAD si benchmark-uri. \\
\addlinespace
VAD &
Detector VAD de baza, folosit ca referinta. &
VAD evaluat separat, etichete adaptive hard/soft, integrare in target-ul de
antrenare si vizualizare live a increderii VAD. \\
\addlinespace
Demo &
Demo offline pentru exemple audio. &
Demo offline plus demo realtime cu microfon, difuzor, surse sintetice, slider SNR,
metode selectabile, snapshot si recording. \\
\addlinespace
Trasabilitate &
Artefacte si rezultate mai greu de separat intre teste si rezultate finale. &
Inventariere automata, clasificare in \texttt{official}, \texttt{benchmark},
\texttt{smoke}, \texttt{demo} si \texttt{legacy}. \\
\bottomrule
\end{longtable}

\section{Imbunatatirea 1: reorganizarea arhitecturii software}

\subsection{Situatia initiala}

In primele iteratii, proiectul functiona ca un ansamblu de scripturi si module
orientate spre demonstrarea ideii principale: reducerea zgomotului cu ajutorul
unui model neuronal si al unor metode clasice. Aceasta abordare este suficienta
pentru prototipare, dar devine limitativa atunci cand proiectul trebuie extins,
evaluat si prezentat intr-o lucrare de licenta.

\subsection{Modificarea introdusa}

Codul a fost reorganizat pe responsabilitati clare:

\begin{itemize}
    \item \texttt{src/data\_prep} pentru pregatirea seturilor de date, generarea
    metadatelor, augmentare si etichete VAD;
    \item \texttt{src/dsp} pentru STFT, metrici, VAD, Wiener Filter si Spectral
    Subtraction;
    \item \texttt{src/models} pentru arhitectura MaskNet si inferenta;
    \item \texttt{src/train} pentru configurarea si rularea antrenarii;
    \item \texttt{src/eval} pentru evaluari si rapoarte;
    \item \texttt{src/tools} pentru CLI, healthcheck, demo realtime si inventar;
    \item \texttt{src/utils} pentru logging, I/O audio, raportare si pipeline-uri.
\end{itemize}

\subsection{Efectul imbunatatirii}

Separarea pe module reduce dependentele implicite si face proiectul mai usor de
inteles. Fiecare componenta poate fi testata si extinsa independent. In plus,
structura actuala permite explicarea mai clara a fluxului complet: date
$\rightarrow$ reprezentare timp-frecventa $\rightarrow$ model $\rightarrow$
reconstructie $\rightarrow$ evaluare.

\section{Imbunatatirea 2: trecerea de la SimpleMaskNet la MaskNet avansat}

\subsection{Situatia initiala}

Modelul initial, pastrat in proiect ca \texttt{SimpleMaskNet}, este o retea
convolutionala simpla. Aceasta contine cateva straturi \texttt{Conv2D},
\texttt{BatchNorm} si activari ReLU, urmate de o convolutie finala si o activare
sigmoid. Scopul sau este estimarea unei masti reale in intervalul $[0,1]$ pentru
spectrograma de magnitudine.

\subsection{Modificarea introdusa}

Modelul principal a fost inlocuit cu o arhitectura MaskNet mai expresiva:

\begin{itemize}
    \item encoder-decoder de tip U-Net cu patru niveluri;
    \item skip connections intre encoder si decoder;
    \item bottleneck convolutional;
    \item BiLSTM pentru modelarea contextului temporal;
    \item mecanisme de channel attention;
    \item variante configurabile: \texttt{tiny}, \texttt{tiny\_fast},
    \texttt{balanced}, \texttt{balanced\_res}, \texttt{large};
    \item blocuri reziduale pentru varianta recomandata \texttt{balanced\_res}.
\end{itemize}

\subsection{Comparatie tehnica}

\begin{center}
\begin{tabular}{lll}
\toprule
\textbf{Caracteristica} & \textbf{SimpleMaskNet} & \textbf{MaskNet actual} \\
\midrule
Tip arhitectura & CNN simplu & U-Net + BiLSTM + Attention \\
Skip connections & Nu & Da \\
Modelare temporala & Limitata & BiLSTM bidirectional \\
Atentie pe canale & Nu & Da \\
Blocuri reziduale & Nu & Optional, recomandat \\
Variante de capacitate & Nu & Da \\
Output & Masca reala & Masca reala sau complexa \\
\bottomrule
\end{tabular}
\end{center}

\subsection{Efectul imbunatatirii}

Noua arhitectura permite modelului sa invete atat structuri locale din
spectrograma, prin convolutii, cat si dependente temporale pe durata mai multor
cadre, prin BiLSTM. Skip connections ajuta la pastrarea detaliilor
timp-frecventa, iar attention-ul permite recalibrarea canalelor relevante.
Blocurile reziduale stabilizeaza antrenarea si imbunatatesc propagarea
gradientilor.

\section{Imbunatatirea 3: introducerea blocurilor reziduale}

\subsection{Situatia initiala}

Arhitectura convolutionala initiala procesa secvential feature-urile, fara cai
explicite de shortcut in interiorul blocurilor. Pe masura ce reteaua devine mai
adanca, aceasta abordare poate duce la optimizare mai dificila si la pierderea
unor informatii utile.

\subsection{Modificarea introdusa}

In varianta \texttt{balanced\_res}, blocurile encoderului, decoderului si
bottleneck-ului pot folosi conexiuni reziduale. Un bloc rezidual poate fi descris
prin:

\begin{equation}
    h = \mathrm{ReLU}(F(x) + P(x)),
\end{equation}

unde $F(x)$ reprezinta transformarile convolutionale, iar $P(x)$ este fie
identitatea, fie o proiectie $1 \times 1$ atunci cand numarul de canale se schimba.

\subsection{Efectul imbunatatirii}

Conexiunile reziduale reduc riscul de degradare la retele mai adanci si permit
pastrarea unor informatii originale in timp ce reteaua invata corectii. Din acest
motiv, presetul recomandat actual este \texttt{stage2\_recommended}, care combina
\texttt{balanced\_res}, AdamW, OneCycleLR si MSE.

\section{Imbunatatirea 4: pipeline de date mai robust}

\subsection{Situatia initiala}

Primele iteratii se bazau pe incarcarea perechilor de fisiere audio si pe
calcularea spectrogramelor necesare pentru antrenare. O astfel de abordare este
functionala, dar vulnerabila la diferente de sample rate, fisiere nepereche,
durate diferite si secvente de lungimi variabile.

\subsection{Modificarea introdusa}

Clasa \texttt{SpeechEnhancementDataset} include acum:

\begin{itemize}
    \item cautare robusta a perechilor clean/noisy;
    \item suport pentru fisiere \texttt{.wav} si \texttt{.flac};
    \item conversie la mono;
    \item resampling automat la 16 kHz;
    \item alinierea semnalelor clean si noisy la aceeasi lungime;
    \item crop optional pentru limitarea duratei;
    \item augmentare audio;
    \item preprocessing;
    \item calcul STFT complex pentru semnalul clean si noisy;
    \item generarea mastii ideale reale si a mastii complexe CRM;
    \item padding in \texttt{collate\_fn\_pad} pentru batch-uri cu lungimi
    diferite.
\end{itemize}

\subsection{Efectul imbunatatirii}

Pipeline-ul actual elimina multe surse de erori practice. Fisierele cu sample
rate diferit pot fi folosite fara conversie manuala, secventele de durate diferite
pot intra in acelasi batch, iar modelul poate fi antrenat atat pe magnitudine, cat
si pe reprezentari complexe.

\section{Imbunatatirea 5: extinderea seturilor de date}

\subsection{Situatia initiala}

In primele forme, proiectul se concentra pe un set de date principal si pe
validarea functionala a metodelor de denoising.

\subsection{Modificarea introdusa}

Forma actuala documenteaza doua configuratii principale:

\begin{itemize}
    \item VoiceBank-DEMAND, cu perechi clean/noisy pentru antrenare si testare;
    \item LibriSpeech + DEMAND, cu fisiere generate la mai multe niveluri SNR
    si cu etichete VAD asociate.
\end{itemize}

Pentru LibriSpeech + DEMAND, pipeline-ul poate resampla automat atat vorbirea
curata, cat si zgomotul, la 16 kHz.

\subsection{Efectul imbunatatirii}

Adaugarea unui al doilea scenariu de date creste relevanta experimentala a
proiectului. VoiceBank-DEMAND ramane util pentru comparatii standard, iar
LibriSpeech + DEMAND permite experimente controlate cu niveluri diferite de zgomot
si integrarea etichetelor VAD.

\section{Imbunatatirea 6: integrarea VAD hard si soft}

\subsection{Situatia initiala}

VAD-ul era prezent ca metoda de detectie a activitatii vocale si ca element de
evaluare separat. Acesta oferea o estimare binara sau simplificata a prezentei
vorbirii.

\subsection{Modificarea introdusa}

Proiectul actual include generarea si utilizarea etichetelor VAD:

\begin{itemize}
    \item etichete hard, pentru decizii binare speech/non-speech;
    \item etichete soft, pentru scoruri continue in intervalul $[0,1]$;
    \item \texttt{vad\_soft\_mask\_floor}, care evita anularea completa a zonelor
    ambigue;
    \item comparatii intre variante fara VAD, cu VAD hard si cu VAD soft.
\end{itemize}

Pentru etichetele soft, gate-ul folosit poate fi scris:

\begin{equation}
    g[t] = f + (1-f)v[t],
\end{equation}

unde $v[t]$ este scorul VAD, iar $f$ este valoarea minima pastrata pentru zonele
incerte.

\subsection{Efectul imbunatatirii}

VAD-ul nu mai este doar o componenta auxiliara, ci devine parte a modului in care
se construieste target-ul de antrenare. In zonele fara vorbire, modelul este
incurajat sa atenueze mai agresiv zgomotul. In zonele ambigue, varianta soft
reduce riscul de stergere a segmentelor slabe de vorbire.

\section{Imbunatatirea 7: introducerea Complex MaskNet}

\subsection{Situatia initiala}

MaskNet clasic estimeaza o masca reala aplicata pe magnitudinea spectrogramei
zgomotoase. Reconstructia audio foloseste magnitudinea estimata si faza semnalului
zgomotos:

\begin{equation}
    \hat{S}[k,t] = \hat{|S|}[k,t] e^{j \angle Y[k,t]}.
\end{equation}

Aceasta metoda este stabila, dar are o limitare importanta: faza zgomotoasa este
refolosita, nu estimata.

\subsection{Modificarea introdusa}

Complex MaskNet foloseste doua canale de intrare, partea reala si partea imaginara
a STFT-ului:

\begin{equation}
    X[k,t] =
    \begin{bmatrix}
        Y_r[k,t] \\
        Y_i[k,t]
    \end{bmatrix}.
\end{equation}

Modelul prezice o masca complexa:

\begin{equation}
    M[k,t] = M_r[k,t] + jM_i[k,t].
\end{equation}

Spectrograma curatata se obtine prin inmultire complexa:

\begin{equation}
    \hat{S}[k,t] = M[k,t]Y[k,t].
\end{equation}

\subsection{Comparatie intre MaskNet clasic si Complex MaskNet}

\begin{center}
\begin{tabular}{lll}
\toprule
\textbf{Aspect} & \textbf{MaskNet clasic} & \textbf{Complex MaskNet} \\
\midrule
Intrare & Magnitudine $|Y|$ & Real si imaginar $(Y_r,Y_i)$ \\
Canale intrare & 1 & 2 \\
Target & Ideal Ratio Mask & Complex Ratio Mask \\
Canale iesire & 1 & 2 \\
Activare iesire & Sigmoid & Tanh scalat \\
Interval output & $[0,1]$ & $[-5,5]$ \\
Faza & Refoloseste faza noisy & Corecteaza implicit faza \\
Reconstructie & Magnitudine + faza noisy & Spectrograma complexa estimata \\
\bottomrule
\end{tabular}
\end{center}

\subsection{Efectul imbunatatirii}

Complex MaskNet face modelul mai expresiv. El poate modifica nu doar energia
spectrala, ci si componenta de faza prin estimarea directa a unei masti complexe.
Aceasta este o imbunatatire importanta fata de abordarea bazata doar pe
magnitudine, mai ales in zonele unde calitatea perceptuala este limitata de faza
semnalului zgomotos.

\section{Imbunatatirea 8: antrenare configurabila prin preseturi}

\subsection{Situatia initiala}

In fazele timpurii, schimbarea configuratiei de antrenare presupunea modificari
manuale sau parametri separati. Acest lucru facea mai dificila reproducerea
experimentelor.

\subsection{Modificarea introdusa}

Configuratia actuala foloseste preseturi:

\begin{itemize}
    \item preseturi de model, precum \texttt{balanced\_res} si
    \texttt{complex\_balanced\_res};
    \item preseturi de optimizer, inclusiv AdamW;
    \item preseturi de scheduler, inclusiv OneCycleLR;
    \item preseturi de loss, precum MSE, MAE, Huber si weighted combined;
    \item preseturi experimentale complete, precum
    \texttt{voicebank\_quality\_final},
    \texttt{librispeech\_soft\_vad\_recommended} si
    \texttt{librispeech\_complex\_masknet\_recommended}.
\end{itemize}

Exemplu de rulare recomandata:

\begin{lstlisting}[language=bash]
python -m src.train.train_mask_model --experiment-preset stage2_recommended --max-batches 30
\end{lstlisting}

\subsection{Efectul imbunatatirii}

Preseturile transforma experimentele in configuratii numite si usor de repetat.
Aceasta este o imbunatatire metodologica importanta: rezultatele pot fi asociate
cu o configuratie clara, iar comparatiile intre variante devin mai corecte.

\section{Imbunatatirea 9: optimizare si stabilitate in training}

\subsection{Situatia initiala}

Primele iteratii validau functionalitatea modelului, dar nu ofereau la fel de
multe mecanisme pentru stabilitatea si monitorizarea antrenarii.

\subsection{Modificarea introdusa}

Forma actuala include:

\begin{itemize}
    \item Adam, AdamW, SGD si RMSprop;
    \item schedulere precum ReduceLROnPlateau, StepLR, CosineAnnealingLR si
    OneCycleLR;
    \item clipping de gradient;
    \item mixed precision;
    \item salvare periodica de checkpoint-uri;
    \item salvarea celui mai bun model pe validare;
    \item resume din checkpoint;
    \item \texttt{training\_history.json}, \texttt{summary.json} si index de
    experimente.
\end{itemize}

\subsection{Efectul imbunatatirii}

Aceste mecanisme reduc riscul de antrenari instabile si permit reluarea
experimentelor. In acelasi timp, istoricul salvat face posibila analiza ulterioara
a evolutiei pierderii, a configuratiei si a performantelor.

\section{Imbunatatirea 10: evaluare comparativa extinsa}

\subsection{Situatia initiala}

Evaluarea initiala era orientata spre verificarea faptului ca metodele produc un
semnal imbunatatit. Aceasta etapa este utila, dar nu este suficienta pentru o
comparatie riguroasa.

\subsection{Modificarea introdusa}

Proiectul actual include scripturi pentru:

\begin{itemize}
    \item evaluarea tuturor metodelor pe acelasi set de fisiere;
    \item compararea noisy baseline, Spectral Subtraction, Wiener Filter si
    MaskNet;
    \item calculul statisticilor si imbunatatirilor;
    \item generarea de tabele si grafice;
    \item compararea checkpoint-urilor;
    \item compararea ultimului experiment cu un baseline;
    \item agregarea benchmark-urilor finale.
\end{itemize}

Exemplu:

\begin{lstlisting}[language=bash]
python -m src.tools.project_cli eval-all -- --clean-dir data/voicebank_demand/clean/test --noisy-dir data/voicebank_demand/noisy/test
\end{lstlisting}

\subsection{Efectul imbunatatirii}

Evaluarea devine reproductibila si comparabila. Metodele clasice nu sunt eliminate,
ci pastrate ca baseline-uri. Astfel, se poate demonstra daca modelul neuronal
aduce un castig real fata de metode DSP traditionale.

\section{Imbunatatirea 11: CLI unificat si healthcheck}

\subsection{Situatia initiala}

Rularea proiectului presupunea cunoasterea scripturilor individuale si a
argumentelor fiecaruia.

\subsection{Modificarea introdusa}

A fost introdus \texttt{project\_cli}, care centralizeaza comenzile principale:

\begin{itemize}
    \item \texttt{train};
    \item \texttt{eval-vad};
    \item \texttt{compare-vad};
    \item \texttt{eval-all};
    \item \texttt{compare-checkpoints};
    \item \texttt{compare-latest};
    \item \texttt{summarize-benchmarks};
    \item \texttt{project-inventory};
    \item \texttt{healthcheck};
    \item \texttt{demo};
    \item \texttt{realtime-demo}.
\end{itemize}

\subsection{Efectul imbunatatirii}

CLI-ul unificat reduce efortul de utilizare si face proiectul mai usor de
prezentat. Healthcheck-ul ajuta la identificarea rapida a problemelor de
configuratie sau de dataset inainte de antrenari lungi.

\section{Imbunatatirea 12: demo offline si realtime}

\subsection{Situatia initiala}

Demo-ul initial permitea rularea pe exemple audio si generarea de rezultate pentru
inspectie.

\subsection{Modificarea introdusa}

Forma actuala include un demo realtime cu:

\begin{itemize}
    \item intrare de la microfon;
    \item iesire catre difuzoare;
    \item metode selectabile: MaskNet, Spectral Subtraction, Wiener Filter si
    bypass;
    \item vizualizare raw vs enhanced;
    \item spectrograme live;
    \item VAD confidence live;
    \item selectie de dispozitive audio;
    \item mod \texttt{speech+noise} cu fisiere incarcate de utilizator;
    \item slider pentru SNR;
    \item scenarii precum \texttt{neutral}, \texttt{office}, \texttt{street} si
    \texttt{hard};
    \item control dry/wet si output gain;
    \item snapshot si inregistrare audio.
\end{itemize}

\subsection{Efectul imbunatatirii}

Demo-ul nu mai este doar o verificare tehnica, ci o componenta de prezentare.
Utilizatorul poate observa diferenta dintre semnalul zgomotos si cel imbunatatit,
poate schimba metoda de denoising si poate testa scenarii diferite fara a reporni
programul.

\section{Imbunatatirea 13: gestionarea artefactelor si a rezultatelor finale}

\subsection{Situatia initiala}

Pe masura ce apar mai multe antrenari, fisiere de rezultate si checkpoint-uri,
devine dificil de stiut care sunt artefactele oficiale si care sunt doar teste.

\subsection{Modificarea introdusa}

Proiectul include un sistem de inventariere care clasifica artefactele in:

\begin{itemize}
    \item \texttt{official};
    \item \texttt{benchmark};
    \item \texttt{smoke};
    \item \texttt{demo};
    \item \texttt{legacy};
    \item \texttt{other}.
\end{itemize}

Comanda dedicata este:

\begin{lstlisting}[language=bash]
python -m src.tools.project_cli project-inventory
\end{lstlisting}

Aceasta genereaza fisiere CSV si JSON cu inventarul artefactelor si cu
recomandarile pentru livrare.

\subsection{Efectul imbunatatirii}

Aceasta etapa imbunatateste calitatea livrarii proiectului. Rezultatele finale pot
fi identificate rapid, iar testele scurte sau rularile istorice raman disponibile
fara sa fie confundate cu artefactele oficiale.

\section{Fluxul actual complet al proiectului}

Fluxul complet in forma actuala poate fi descris astfel:

\begin{enumerate}
    \item Se verifica structura datasetului cu \texttt{healthcheck}.
    \item Se pregatesc datele clean/noisy si, daca este cazul, etichetele VAD.
    \item Se selecteaza un preset experimental.
    \item Dataloader-ul incarca perechi clean/noisy, resampleaza, aplica
    preprocessing si calculeaza STFT.
    \item Pentru MaskNet clasic se construieste IRM; pentru Complex MaskNet se
    construieste CRM.
    \item Modelul este antrenat cu optimizerul si schedulerul configurate.
    \item Se salveaza checkpoint-uri si istoricul de antrenare.
    \item Se evalueaza checkpoint-ul fata de baseline-uri clasice.
    \item Se genereaza rapoarte, grafice si tabele.
    \item Se ruleaza demo-ul offline sau realtime.
    \item Se inventariaza artefactele finale pentru lucrare si prezentare.
\end{enumerate}

\section{Comenzi reprezentative}

\subsection{Verificarea proiectului}

\begin{lstlisting}[language=bash]
python -m src.tools.project_cli healthcheck -- --dataset librispeech
\end{lstlisting}

\subsection{Antrenare recomandata pe LibriSpeech}

\begin{lstlisting}[language=bash]
python -m src.train.train_mask_model --experiment-preset librispeech_soft_vad_recommended
\end{lstlisting}

\subsection{Antrenare recomandata cu VAD soft pe LibriSpeech}

\begin{lstlisting}[language=bash]
python -m src.data_prep.generate_vad_labels --split all --label-set adaptive_soft_v1 --label-mode soft
python -m src.train.train_mask_model --experiment-preset librispeech_soft_vad_recommended
\end{lstlisting}

\subsection{Antrenare Complex MaskNet}

\begin{lstlisting}[language=bash]
python -m src.train.train_mask_model --experiment-preset librispeech_complex_masknet_recommended
\end{lstlisting}

\subsection{Compararea checkpoint-urilor}

\begin{lstlisting}[language=bash]
python -m src.tools.project_cli compare-checkpoints -- --clean-dir data/voicebank_demand/clean/test --noisy-dir data/voicebank_demand/noisy/test --checkpoint baseline=checkpoints/baseline.pth --checkpoint stage2=checkpoints/stage2_recommended.pth --max-files 20
\end{lstlisting}

\subsection{Demo realtime}

\begin{lstlisting}[language=bash]
python -m src.tools.project_cli realtime-demo -- --method masknet --visualize --record
\end{lstlisting}

\section{Concluzie}

Comparativ cu primele iteratii, forma actuala a proiectului este substantial mai
completa. Imbunatatirile nu constau doar in cresterea complexitatii modelului, ci
si in maturizarea intregului ecosistem experimental: date mai bine gestionate,
antrenare reproductibila, evaluare comparativa, demo in timp real si organizare
clara a artefactelor.

Din punct de vedere tehnic, evolutia principala este trecerea de la o retea simpla
de estimare a unei masti reale la o familie de modele MaskNet configurabile,
incluzand varianta reziduala recomandata si varianta Complex MaskNet. Din punct de
vedere metodologic, proiectul a evoluat de la un prototip functional la o platforma
care poate sustine experimente, comparatii si demonstratii reproductibile.

Prin aceste imbunatatiri, proiectul devine mai potrivit pentru o lucrare de
licenta: are o structura explicabila, rezultate comparabile, componente
demonstrabile si o separare clara intre prototipuri, teste si artefacte finale.

\end{document}
